% Inactive proof reconstruction. The compact Hermitian setting of Chapter12
% and 2<=m<n are in force. No external bounded comparison theorem is used.
% Published Gu--Nguyen Lemmas3.8--3.10 supply a related local mechanism;
% their bounded-product hypotheses are not silently transferred here.

For bounded admissible $u,V$, use the notation
\[
 \sigma_u=\omega+\ddc u,\qquad
 M_{p,q}=\sigma_u^p\sigma_V^q\omega^{n-p-q}
                    \quad(0\leq p+q\leq m).
\]
These expressions denote the positive mixed measure $Q_m$ with
$p$ slots equal to $u$, $q$ slots equal to $V$ and the remaining
slots equal to the constant zero. In particular $M_{0,0}=\omega^n$.
For $p+q\leq m-1$, put
\[
 T_{p,q}=\sigma_u^p\sigma_V^q\omega^{n-p-q-1}.
\]
This denotes $R_{m-1}$ with the same nonconstant slots and zero
in its other slots, wedged with $\omega^{n-m}$.

\begin{lemma}\label{draft637:mixed-comparison}
Let $u,V$ be bounded admissible, $b$ a real constant and $s>0$.
Suppose $0\leq\max\{u,V+b\}-u\leq s$. Set
$E=\{u<V+b\}$ and $b_{p,q}=M_{p,q}(E)$.
There is a constant $C$, depending only on $X,\omega,m,n$, such that
\[
 b_{p,q+1}\leq b_{p+1,q}
              +Cs\sum_{a+d\leq p+q}b_{a,d}
                   \qquad(p+q\leq m-1).
\]
For $q=0$, the error can instead be bounded by
$Cs(b_{p,0}+b_{p-1,0}+b_{p-2,0})$, with negative indices zero.
\end{lemma}
\begin{proof}
First estimate the torsion of $T_{p,q}$, writing $r=p+q$.
For smooth admissible potentials, the product rule expresses
$\ddc T_{p,q}$ as a sum of fixed smooth metric forms wedged
with products of $\sigma_u,\sigma_V$ of degrees $r,r-1,r-2$.
The terms of degree $r$ come from differentiating the background
power; they cannot be discarded. Indeed,
$\mathrm d\sigma_u=\mathrm d\sigma_V=\mathrm d\omega$,
$\dc\sigma_u=\dc\sigma_V=\dc\omega$ and
$\ddc\sigma_u=\ddc\sigma_V=\ddc\omega$.
Every term therefore contains either $\ddc\omega$ or
$\mathrm d\omega\wedge\dc\omega$, with its product-rule
multiplicity and the appropriate remaining background power.
Terms with negative numbers of factors are absent.

All variable degrees are at most $r<m$.
Apply \cref{hessian:absolute-mixed-forms} to each term. Expanding
the resulting positive power of the sum of its variable factors
gives the pointwise bound
\[
 |\ddc T_{p,q}|\leq C\sum_{a+d\leq r}M_{a,d}.
\]
For $q=0$ every variable factor is $\sigma_u$, so the sharper
bound is $C(M_{p,0}+M_{p-1,0}+M_{p-2,0})$.
The constant bounds fixed metric coefficients on a finite coordinate
cover and the finitely many multiplicities; it is independent of
the potentials and their bounds.

Choose uniformly bounded smooth decreasing approximations to $u,V$.
By \cref{hessian:bounded-torsion-weight}, the torsion measures
converge weakly; by \cref{hessian:joint-decreasing-products},
all positive measures in these bounds also converge weakly.
Against nonnegative continuous tests, lower semicontinuity of
total variation preserves both bounds. They are consequently
inequalities of measures and may be restricted to the Borel set $E$.

Put $w=\max\{u,V+b\}$ and $h=w-u$.
The mixed maximum theorem \cref{hessian:bounded-maximum-measure},
with one moving slot and the fixed slots defining $T_{p,q}$, gives
\[
 (\omega+\ddc w)T_{p,q}
 \geq\mathbf1_E\sigma_VT_{p,q}
                +\mathbf1_{X\setminus E}\sigma_uT_{p,q}.
\]
All products here are the corresponding positive mixed measures.
After integration, this implies
\[
 b_{p,q+1}\leq b_{p+1,q}
       +\int_X\bigl((\omega+\ddc w)-\sigma_u\bigr)T_{p,q}.
\]
Regularized maxima of the smooth decreasing approximations give
a uniformly bounded decreasing smooth list $w_j\downarrow w$.
Stokes' theorem on compact $X$ gives for these lists
\[
 \int_X\bigl((\omega+\ddc w_j)-\sigma_{u_j}\bigr)T_{p,q}^j
                      =\int_X(w_j-u_j)\ddc T_{p,q}^j.
\]
Joint decreasing convergence passes the positive mixed measures
on the left to their limits. The continuous two-variable weight
$F(x,y)=y-x$ in \cref{hessian:bounded-torsion-weight} passes the
right side to $\int_Xh\ddc T_{p,q}$. No monotonicity of $w_j-u_j$
is needed. Since $h$ is zero outside $E$ and $0\leq h\leq s$,
\[
 \int_Xh\ddc T_{p,q}\leq s|\ddc T_{p,q}|(E).
\]
The torsion bounds prove both assertions. The argument never
passes a varying measure to a weak limit after restriction to $E$.
\end{proof}

\begin{theorem}\label{draft637:modified-comparison}
There is $C\geq1$, depending only on $X,\omega,m,n$, with the
following property. For bounded admissible $u,v$ and $0<\alpha<1$,
put $\eta=1-\alpha$, $V=\alpha v$ and $S=\inf_X(u-V)$.
For $0<s<\eta^3/C$, let $E_s=\{u<V+S+s\}$. Then
\[
 \int_{E_s}H_m(V)
       \leq\left(1+\frac{Cs}{\eta^m}\right)
                                  \int_{E_s}H_m(u).
\]
No differential identity for a power of $\omega$ is assumed.
\end{theorem}
\begin{proof}
Use the mixed measures above on $E_s$, writing
$b_{p,q}=M_{p,q}(E_s)$ and $a_k=b_{k,0}$.
The definition of $S$ gives
$0\leq\max\{u,V+S+s\}-u\leq s$, so the preceding lemma applies.
Repeatedly moving a $V$ slot to a $u$ slot gives, for $p+q=r$,
\[
 b_{p,q}\leq a_r+C_1s\sum_{d=0}^{r-1}A_d,
 \qquad A_d=\max_{a+b=d}b_{a,b},\qquad 1\leq r\leq m.
\]
Here there are at most $m$ moves, and the multiplicity $d+1$
in the lower-degree sum is absorbed in $C_1$.
For $0<s\leq1$, induction on $r$ shows
\[
 A_r\leq a_r+K_rs\sum_{k<r}a_k.
\]
In fact $A_0=a_0$, and one may take
$K_r=C_1(1+\sum_{1\leq d<r}K_d)$: insert the earlier bounds
in the displayed recurrence and use $s^2\leq s$.
Thus, for a constant $C_2$ depending only on the setting,
\[
 b_{0,m}\leq a_m+C_2s\sum_{k<m}a_k.
\]

It remains to control the lower diagonal masses by $a_m$.
The identity $\sigma_V=\alpha\sigma_v+\eta\omega$ and mixed
positivity give $b_{k,1}\geq\eta a_k$; this identity and
inequality pass from smooth approximations by joint convergence.
The sharper diagonal assertion of the lemma therefore gives
\[
 \eta a_k\leq a_{k+1}
          +C_3s(a_k+a_{k-1}+a_{k-2})\qquad(0\leq k<m),
\]
with $a_{-1}=a_{-2}=0$.
Set $B_k=\eta^{m-k}a_k$ for $0\leq k\leq m$, and set
$B_{-1}=B_{-2}=0$. Multiplying by $\eta^{m-k-1}$ and using
$0<\eta<1$ gives
\[
 B_k\leq B_{k+1}
          +\delta(B_k+B_{k-1}+B_{k-2}),
 \qquad\delta=C_3s\eta^{-3}.
\]
For $M=\max_{0\leq k<m}B_k$, summing from any $j<m$ through
$m-1$ gives $B_j\leq B_m+3m\delta M$.
If $s<\eta^3/(6mC_3)$, it follows that $M\leq2B_m$.
Consequently
\[
 a_k\leq2\eta^{-(m-k)}a_m\quad(k<m),\qquad
 \sum_{k<m}a_k\leq2m\eta^{-m}a_m.
\]
Choose $C\geq\max\{1,6mC_3,2mC_2\}$. Then the asserted
range for $s$ supplies every preceding restriction, and substitution
in the bound for $b_{0,m}$ proves the theorem.
No division by $a_m$ occurs, so zero top mass on $E_s$ is included.
\end{proof}

\begin{theorem}\label{draft637:strict-domination}
Let $u,v$ be bounded admissible. If, for some $0\leq c<1$,
\[
 \mathbf1_{\{u<v\}}H_m(u)
                  \leq c\mathbf1_{\{u<v\}}H_m(v),
\]
then $u\geq v$ everywhere, on an arbitrary Hermitian background.
In particular, $H_m(u)(\{u<v\})=0$ implies $u\geq v$.
\end{theorem}
\begin{proof}
Add the same constant to $u,v$ so that $v\geq1$.
This preserves the measure hypothesis and the desired conclusion.
Choose $0<\alpha<1$ with $\alpha^m>c$ and put $V=\alpha v$,
$\eta=1-\alpha$. Suppose $S=\inf_X(u-V)<0$.
For $0<s<-S$, the set $E_s=\{u<V+S+s\}$ is contained in
$\{u<V\}$ and hence in $\{u<v\}$.
It has positive volume. Indeed, if its volume were zero, then
$u\geq V+S+s$ almost everywhere. Both functions are bounded
admissible with their canonical representatives, so
\cref{hessian:canonical-order} would give this order everywhere,
contradicting the definition of $S$. The argument does not require
$E_s$ to be open or the infimum to be attained.

Choose also $s<\eta^3/C$, where $C$ is from the preceding theorem.
By \cref{hessian:convex-background-term} and the measure hypothesis,
\[
 \alpha^m H_m(v)(E_s)+\eta^m\int_{E_s}\omega^n
 \leq H_m(V)(E_s)
 \leq(1+Cs\eta^{-m})H_m(u)(E_s)
 \leq c(1+Cs\eta^{-m})H_m(v)(E_s).
\]
If $c>0$, take $s$ still smaller so that
$c(1+Cs\eta^{-m})<\alpha^m$. If $c=0$, this strict inequality
already holds. All measures are finite and positive, so the
display forces the positive-volume term to be zero, a contradiction.
Hence $S\geq0$ and $u\geq\alpha v$. Let $\alpha$ increase
to one, with $\alpha^m>c$. Since $v$ is bounded, $\alpha v\to v$
uniformly, and therefore $u\geq v$. The final assertion is $c=0$.
\end{proof}

% If integrated, the existing proper bounded uniqueness, data stability and
% smoothness proofs use only strict domination after their hypotheses.
% Their partial-power assumptions can then be removed after checking callers.
% This does not assert existence for arbitrary Radon data, c=1 domination,
% unmodified global comparison, a uniform mass lower bound, or novelty.
